% profiles/the-university-of-texas-at-austin.tex
% — The University of Texas at Austin
% ============================================================================
% source: https://graduate.utexas.edu/navigating/graduation/digital-submission
%         (Office of Graduate and Postdoctoral Studies, Digital Submission
%          Requirement — the page that links the guidelines below)
% source: "Format Guidelines.pdf" — "Graduate School Format Guidelines for
%         Dissertations, Treatises, Theses, and Reports", 18 pp., internal PDF
%         title "Graduate School Format Guidelines - Proposed Revisions[2]",
%         embedded ModDate 2025-02-14, linked from that page as "Format
%         Guidelines":
%         https://utexas.box.com/s/ffag6al7b0eaaqcheml6dcot9u3z6zit
% source: https://graduate.utexas.edu/navigating/graduation/digital-submission/latex
%         and the "LaTex Dissertation/Thesis package (ZIP Format)" it links,
%         https://utexas.box.com/s/tfzk7cjur5zbnblnpcf1l8m884t03786
%         (UT's OFFICIAL LaTeX package; the page says "The styling contained in
%          these packages meets the requirements of the Graduate School as of
%          April 2023". Read for the margin note below — NOT used to set one.)
% accessed: 2026-09-05
% checked:  texlib-thesis-profile-round
%
% VERIFIED and set here: spacing, the title page, and the committee membership
% page. NOT SET, deliberately: \thessissetgeometry — see immediately below. It
% is the important part of this file.
%
% Box access recipe, same as the Washington University profile: the /s/<hash>
% pages are JavaScript viewers curl cannot read. Open in a browser, take the
% f_<digits> file id out of the DOM, then fetch
%   https://utexas.app.box.com/index.php?rm=box_download_shared_file&shared_name=<hash>&file_id=f_<digits>
%
% ============================================================================
% UT AUSTIN PUBLISHES NO MARGIN FIGURE. THE GEOMETRY IS UNSET ON PURPOSE.
%
% The Format Guidelines' entire margin section, quoted in full, is:
%
%   "b. Margins
%    Margins should be consistent throughout the document, including the
%    appendix."
%
% That is a CONSISTENCY rule, not a measurement. It was checked exhaustively:
% across all 18 pages the word "inch" occurs exactly once, and it is about
% oversized plates ("more than 11x14 inches"); no figure with an inch mark
% appears anywhere near the margin section. So there is no published number to
% record, and under RESEARCHING.md an unverified figure stays unset. The class's
% neutral geometry renders, which satisfies UT's only stated rule (it is
% consistent) while making no claim about a figure UT did not publish.
%
% WHAT UT'S OWN LaTeX PACKAGE USES, recorded because it is genuinely useful and
% NOT because it is a requirement. styles/utformat.sty sets:
%
%   \oddsidemargin 0.25in  \evensidemargin 0.25in   -> 1.25in left and right
%   \textwidth 6in            % its own comment: "8.5 inches, with 1.25 inch
%                             %  margins = 8.5 - 2 * 1.25 = 6"
%   \topmargin 0.25in                                -> 1.25in top
%   \textheight 8.5in                                -> 1.25in bottom
%
% i.e. 1.25in on all four sides. A filer who wants to match what the Graduate
% School's own package produces can write, in the document preamble:
%
%   \thesissetgeometry{left=1.25in, right=1.25in, top=1.25in, bottom=1.25in}
%
% This profile does NOT do that for them. Setting it would state as UT's
% requirement a figure UT declined to state, and 1in is equally in spec because
% the only rule is consistency. The evidence is laid out here so the choice is
% the filer's and the reviewer's, not this file's.
% ============================================================================

\thesisinstitution{The University of Texas at Austin}

% --- Filing figures ----------------------------------------------------------
% \thesissetgeometry is deliberately absent -- see above.

% "Spacing should be consistent throughout the text. 1.5 or double spacing is
% recommended for ease of reading."
%
% RECOMMENDED, not required, and this line is therefore a default rather than a
% rule -- the same shape as the margin finding, one step less severe. Double is
% one of the two values UT names and is the class default; onehalf is equally
% within the recommendation. The call is made explicitly so that a later change
% to the class's neutral default cannot silently move a UT filing off one of the
% two recommended values.
\thesissetspacing{double}

% No chapter-opening rule beyond "Each new chapter or major section (i.e.,
% Chapter 1, Chapter 2, Appendix, Bibliography, Vita) must begin on a new page",
% which the class already does. No deeper top margin is specified, so
% \thesissetchapteropening stays unset.

% --- Accessibility -----------------------------------------------------------
% NOTHING IS DECLARED REQUIRED, and at UT Austin that is worth a second look
% rather than a shrug. The 18-page Format Guidelines contain ZERO occurrences of
% "accessib", "alt text", "WCAG", "tagged" or "PDF/A"; the Digital Submission
% Requirement page has none; the LaTeX & Overleaf page has none.
%
% UT Austin does publish a university-wide Digital Accessibility Policy
% (utexas.edu/digital-accessibility-policy, linked from the footer of every page
% including these), and as a large public entity it falls under the April 2026
% ADA Title II compliance date. But the Graduate School has not applied either to
% a filed ETD in its published filing requirements, and this class must not
% invent the bridge. Compare the University of Florida profile, where the
% graduate school DID make that link explicitly and the flags are raised.
\thesisaccessibilityrequirement{Not published}
	{The UT Austin Graduate School publishes no accessibility requirement for a
	 filed ETD in its Format Guidelines, its Digital Submission Requirement page
	 or its LaTeX page, as read on 2026-09-05. The university does have a
	 separate, site-wide Digital Accessibility Policy that the Graduate School
	 does not cite in its filing requirements. A finding, not a licence, and one
	 to re-check: UT Austin is a large public entity under the April 2026 ADA
	 Title II date.}

% --- Title page --------------------------------------------------------------
% Reproduced from Sample D (Format Guidelines p. 16, rendered at 85 dpi and read
% as an image). Everything is centred.
%
% ON THE BOLD IN THE SPECIMEN: the sample sets the title, the student name, the
% document type and the degree in bold, and sets nothing else in bold. Those are
% exactly the four placeholders a student replaces, so the bold marks FILL-IN
% FIELDS, not a typographic requirement -- the surrounding prose asks for bold
% nowhere and says only that headings "may be bolded". Everything here is
% therefore plain. Flagged because it is the kind of thing a later pass might
% "correct" from the picture.
%
% The title is set double-spaced: the specimen's own placeholder text reads
% "Title of ETD Centered / and Double-Spaced", which is an instruction wearing
% the costume of a title.
%
% "Month (May, August, or December only) Year" -- so \graddate{May 2027}, and
% only those three months.
%
% "The degree sought must be worded in the form given in the Graduate Catalog.
% If dual degree, list both degrees on a separate line separated by the word
% And." A document supplies that through \graddegree.
%
% The parenthetical lines on the specimen ("(The degree sought must be worded
% ...)") are instructions to the student and are not part of the page.
\renewcommand{\thesistitlepage}{%
	\loadgeometry{nopagenumber}%
	\begin{titlepage}%
		\thispagestyle{empty}\singlespacing\centering
		\vspace*{\fill}
		\begingroup\doublespacing
		\thesis@tagtitle{Title}{\@title}\par
		\endgroup
		\vspace{2\baselineskip}
		by\par
		\vspace{\baselineskip}
		\@author\par
		\vspace{2\baselineskip}
		\thesis@DocNoun\par
		\vspace{\baselineskip}
		Presented to the Faculty of the Graduate School\\
		of the University of Texas at Austin\\
		in Partial Fulfillment\\
		of the Requirements\\
		for the Degree(s) of\par
		\vspace{\baselineskip}
		\thesis@degree\par
		\vspace{2\baselineskip}
		The University of Texas at Austin\\
		\thesis@date\par
		\vspace*{\fill}
	\end{titlepage}%
	\loadgeometry{normal}%
}

% --- Committee membership page ---------------------------------------------------
% UT calls this the COMMITTEE MEMBERSHIP PAGE (also "Committee Listing Page"),
% and it IS part of the filed document -- it is item 2 in the required order,
% ahead of the title page. Reproduced from Sample C (Format Guidelines p. 15).
%
% NO SIGNATURES, AND THAT IS EXPLICIT: "The committee membership page in the pdf
% file that is uploaded for archiving and publication must not contain committee
% signatures." So no rules are drawn -- typed names only. This is the strongest
% such statement of any institution in this round.
%
% PHYSICAL SIGNATURES ARE STILL COLLECTED, on a separate form, and UT is unusually
% specific about them: doctoral students submit the "report of dissertation
% committee" signed by ALL committee members and a GSC representative, "no proxy
% signatures allowed"; "Signatures should all be on a single page"; "Scanned or
% electronic signatures will be accepted as long as they are legible and dark
% enough to be imaged"; "Typed names as a signature are not allowed"; and
% "Extensions will not be granted because a committee member was not available
% to sign."
%
% Layout, from the specimen: a centred certification sentence naming the author;
% the title, centred and double-spaced; then a RIGHT-ALIGNED block headed
% "Committee:" with one name per line.
%
% "Educational or professional titles (Ph.D. or Dr.) are not included for
% committee members, but the titles indicating committee supervisory status
% ('Supervisor' or 'Co-supervisor') must follow the names of committee members
% who have been officially designated in these roles." So \committeemember takes
% a bare name and, where applicable, exactly "Supervisor" or "Co-Supervisor" --
% and NO post-nominals, the same trap as Washington University in St. Louis and
% the opposite of most other profiles here.
%
% A NUMBERING AMBIGUITY, left rather than resolved: UT says "Pages should be
% numbered sequentially starting with 1 on the first page of the document (the
% Copyright Page or Committee Listing Page)" and that numbers are "centered at
% the bottom of the page throughout the document" -- but Samples C and D show no
% page number on either page. This profile follows the SAMPLES and leaves both
% unnumbered, reading the prose as fixing where the COUNT starts. A reviewer who
% reads it as "print 1 on the committee page" should say so.
\renewcommand{\thesisapprovalpage}{%
	\clearpage\thispagestyle{empty}%
	\loadgeometry{nopagenumber}\singlespacing
	\thesis@tagtitle{H1}{{\color{white}\tiny Committee Membership Page}}\par
	\begin{center}
		\vspace*{0.5in}
		The \thesis@DocNoun{} Committee for \@author\\
		certifies that this is the approved version of the following
		\thesis@docnoun:\par
		\vspace{2\baselineskip}
		\begingroup\doublespacing
		\@title\par
		\endgroup
	\end{center}
	\vspace{2\baselineskip}
	% Right-aligned committee block, per the specimen.
	\begingroup
	\raggedleft
	\renewcommand{\thesis@memberline}[2]{%
		\vskip 0.8\baselineskip
		##1\ifx\relax##2\relax\else, ##2\fi\par
	}%
	\textbf{Committee}:\par
	\thesis@committee
	\par
	\endgroup
	\loadgeometry{normal}\doublespacing\clearpage
}
